\documentclass[11pt]{article}
\usepackage[margin=1in]{geometry}
\usepackage{amsmath,amssymb}
\usepackage[T1]{fontenc}
\usepackage{lmodern}
\usepackage[hidelinks]{hyperref}

\newcommand{\F}{\mathcal F}
\newcommand{\R}{\mathbb R}
\newcommand{\MAP}{\mathrm{MAP}}

\title{Literature-Implied One-Dimensional Subcase for arXiv:1501.07036}
\author{fa\_banach\_001 agent note}
\date{2026-06-19}

\begin{document}
\maketitle

\section*{Status}

\textbf{Status:} literature-implied answer, partial subcase.

The source paper is E. Pernecka and R. J. Smith,
\emph{The metric approximation property and Lipschitz-free spaces over subsets
of \(\R^N\)}, arXiv:1501.07036.  After Corollary 1.2, the authors ask whether
\(\F(M)\) has the metric approximation property for all subsets
\(M\subset \R^N\).

The supporting paper is A. Godard, \emph{Tree metrics and their
Lipschitz-free spaces}, arXiv:0904.3178.

\section*{Recorded Subcase}

For every subset \(M\subset \R\), equipped with the inherited Euclidean
metric, the Lipschitz-free space \(\F(M)\) has the metric approximation
property.

More generally, if \(A\) is a subset of an \(\R\)-tree \(T\) such that
\(\operatorname{Br}(T)\subset \overline A\), then \(\F(A)\) has the metric
approximation property.

\section*{Identification}

Godard's Theorem 3.2 states that, for an \(\R\)-tree \(T\) and a subset
\(A\subset T\) with \(\operatorname{Br}(T)\subset \overline A\), the free space
\(\F(A)\) is isometric to \(L_1(\mu_{\overline A})\), where
\(\mu_{\overline A}\) is the explicit length-plus-gap measure associated with
\(\overline A\).  In the special case \(T=\R\), the tree has no branching
points, so the hypothesis is automatic for every subset \(M\subset\R\).
Therefore \(\F(M)\) is isometric to an \(L_1\) space.

Every \(L_1(\mu)\) space has \(\MAP\): finite measurable partitions yield
norm-one conditional expectations with finite-dimensional ranges, and these
expectations converge strongly to the identity along the directed set of
finite subalgebras.  Hence \(\F(M)\) has \(\MAP\) for all \(M\subset\R\).

\section*{Scope}

This is not a full solution of the Pernecka--Smith question.  It covers the
complete one-dimensional case \(N=1\), and the same reasoning covers the
tree-metric subcase above.  The arbitrary subset problem in \(\R^N\) for
\(N\ge 2\) remains open in this packet.

\section*{References}

\begin{itemize}
\item E. Pernecka and R. J. Smith, \emph{The metric approximation property and
Lipschitz-free spaces over subsets of \(\R^N\)}, arXiv:1501.07036.
\item A. Godard, \emph{Tree metrics and their Lipschitz-free spaces},
arXiv:0904.3178.
\end{itemize}

\end{document}
